% ============================================================================
% CAPITULO 5: GEOMETRIA COMPUTACIONAL
% ============================================================================
\section{Geometria Computacional}

% ----------------------------------------------------------------------------
% SUBTEMA 5.1: PRIMITIVAS 2D, TRIGONOMETRIA E UNIAO DE INTERVALOS 1D
% ----------------------------------------------------------------------------
\subsection{Primitivas 2D, Trigonometria e União de Intervalos 1D}

\subsubsection{Struct \texttt{Pt<T>} Completa, Segmentos, Retas e Polígonos}
\algobox
  {Base universal para geometria 2D inteira (\texttt{ll} / \texttt{i128}) e contínua (\texttt{ld}).}
  {Evitar \texttt{double} e funções trigonométricas sempre que possível: usar produto vetorial (\texttt{cross}) e escalar (\texttt{dot}) em inteiros de 64/128 bits. \texttt{cross(a, b, c) > 0} indica curva à esquerda (anti-horário / CCW).}
  {Para coordenadas até $10^9$, \texttt{cross} e \texttt{dist2} podem atingir $4 \cdot 10^{18}$ (cabe no limite de \texttt{ll}, mas em comparações de frações ou circuncentros use \texttt{i128}!). Em ponto flutuante, usar \texttt{EPS = 1e-9} e \texttt{atan2l(y, x)}.}
  {Operações Básicas: $\Ocomp{1}$ $\mid$ Área de Polígono: $\Ocomp{N}$}

\begin{lstlisting}
constexpr ld EPS = 1e-9L;
constexpr ld PI = 3.141592653589793238462643383279502884L;

template <class T> int sgn(T x) { return (x > 0) - (x < 0); }
template <> int sgn(ld x) { return (x > EPS) - (x < -EPS); }

template <class T>
struct Pt {
  T x = 0, y = 0;
  Pt operator+(const Pt& o) const { return {x + o.x, y + o.y}; }
  Pt operator-(const Pt& o) const { return {x - o.x, y - o.y}; }
  Pt operator*(T k) const { return {x * k, y * k}; }
  Pt operator/(T k) const { return {x / k, y / k}; }
  bool operator==(const Pt& o) const { return sgn(x - o.x) == 0 && sgn(y - o.y) == 0; }
  bool operator<(const Pt& o) const {
    return sgn(x - o.x) != 0 ? sgn(x - o.x) < 0 : sgn(y - o.y) < 0;
  }
  T dot(const Pt& o) const { return x * o.x + y * o.y; }
  T cross(const Pt& o) const { return x * o.y - y * o.x; }
  T cross(const Pt& a, const Pt& b) const { return (a - *this).cross(b - *this); }
  T dist2() const { return x * x + y * y; }
  ld dist() const { return hypotl((ld)x, (ld)y); }
  ld angle() const { return atan2l((ld)y, (ld)x); } /* [-pi, pi] */
  Pt perp() const { return {-y, x}; } /* Rotacao +90 graus (CCW) */
  Pt rot(ld a) const {
    ld c = cosl(a), s = sinl(a);
    return {(T)(x * c - y * s), (T)(x * s + y * c)};
  }
};
using Pll = Pt<ll>;
using Pld = Pt<ld>;

/* Orientacao de (a -> b -> c): +1 = CCW (esq), -1 = CW (dir), 0 = colinear */
template <class T> int ccw(const Pt<T>& a, const Pt<T>& b, const Pt<T>& c) {
  return sgn(a.cross(b, c));
}

/* Verifica se o ponto p pertence ao segmento fechado [a, b] */
template <class T> bool on_segment(const Pt<T>& a, const Pt<T>& b, const Pt<T>& p) {
  return ccw(a, b, p) == 0 && sgn((a - p).dot(b - p)) <= 0;
}

/* Intersecao de segmentos fechados [a, b] e [c, d] (exata para inteiros) */
template <class T> bool seg_intersect(const Pt<T>& a, const Pt<T>& b,
                                      const Pt<T>& c, const Pt<T>& d) {
  int oa = ccw(c, d, a), ob = ccw(c, d, b), oc = ccw(a, b, c), od = ccw(a, b, d);
  if (oa != ob && oc != od) return true;
  return on_segment(c, d, a) || on_segment(c, d, b) ||
         on_segment(a, b, c) || on_segment(a, b, d);
}

/* Intersecao de retas infinitas (a + t*v) e (c + u*w) em ponto flutuante */
bool line_intersect(Pld a, Pld b, Pld c, Pld d, Pld& out) {
  Pld v = b - a, w = d - c;
  ld cr = v.cross(w);
  if (sgn(cr) == 0) return false;
  out = a + v * ((c - a).cross(w) / cr);
  return true;
}

/* Dobro da area com sinal de um poligono (positivo se CCW) */
ll polygon_area2(const vector<Pll>& p) {
  ll a = 0; int n = sz(p);
  rep(i, 0, n) a += p[i].cross(p[(i + 1) % n]);
  return a;
}
\end{lstlisting}

\subsubsection{União de Intervalos 1D e Projeção de Sombras}
\algobox
  {Calcular a medida da união de $N$ intervalos $[L_i, R_i]$ na reta (ex.: sombras $[X_i, X_i + H_i \cot \theta]$ projetadas por raios com ângulo $\theta$ em graus).}
  {Ordenar os intervalos por $L_i$ crescente e fundir todo intervalo que intersecta o intervalo acumulado atual $[curL, curR]$. Como $\cot \theta = \frac{1}{\tan(\theta \cdot \pi / 180)}$ é constante para todos os objetos, podemos fazer a união diretamente em ponto flutuante (ou escalar) em $\Ocomp{N \log N}$.}
  {Converter graus para radianos com \texttt{rad = theta * PI / 180.0L}.}
  {Tempo: $\Ocomp{N \log N}$ $\mid$ Espaço: $\Ocomp{N}$}

\begin{lstlisting}
ld union_of_intervals_length(vector<pair<ld, ld>> segs) {
  if (segs.empty()) return 0.0L;
  sort(all(segs));
  ld tot = 0, curL = segs[0].fst, curR = segs[0].snd;
  for (auto [l, r] : segs) {
    if (l > curR) { tot += curR - curL; curL = l; curR = r; }
    else chmax(curR, r);
  }
  return tot + (curR - curL);
}

ld solve_jigsaw_of_shadows(int theta_deg, const vector<pll>& flatlanders) {
  ld cot_t = 1.0L / tanl(theta_deg * PI / 180.0L);
  vector<pair<ld, ld>> segs;
  segs.reserve(sz(flatlanders));
  for (auto [x, h] : flatlanders) segs.eb((ld)x, (ld)x + (ld)h * cot_t);
  return union_of_intervals_length(move(segs));
}
\end{lstlisting}

% ----------------------------------------------------------------------------
% SUBTEMA 5.2: MANHATTAN 45 GRAUS, PARABOLAS E RETAS (PBS + CHT)
% ----------------------------------------------------------------------------
\subsection{Geometria Analítica: Manhattan $45^\circ$, Parábolas e Retas}

\subsubsection{Rotação de Coordenadas de $45^\circ$ (Distância de Manhattan)}
\algobox
  {Transformar problemas de Distância de Manhattan $|x_1 - x_2| + |y_1 - y_2|$ em Distância de Chebyshev $\max(|u_1 - u_2|, |v_1 - v_2|)$, e encontrar todos os pontos inteiros que estão a distância exata $D_i$ de cada torre $(X_i, Y_i)$.}
  {Aplicando a bijeção $u = x + y$ e $v = x - y$:
  $\bullet$ $|x_1 - x_2| + |y_1 - y_2| = \max(|u_1 - u_2|, |v_1 - v_2|)$.
  $\bullet$ Um par $(u, v) \in \ZZ^2$ corresponde a um ponto inteiro $(x, y) = (\frac{u+v}{2}, \frac{u-v}{2}) \in \ZZ^2$ sse $u \equiv v \pmod 2$ (mesma paridade!).
  $\bullet$ O conjunto de pontos a distância de Manhattan \textbf{exatamente} $D_i$ de $(X_i, Y_i)$ (uma "esfera de Manhattan") vira a \textbf{borda de um quadrado alinhado aos eixos} $[U_i - D_i, U_i + D_i] \times [V_i - D_i, V_i + D_i]$ no plano $(u, v)$, que consiste em apenas 4 segmentos ortogonais (2 horizontais e 2 verticais)!
  Para achar a interseção das $N$ bordas de quadrados: começamos com os 4 segmentos da primeira torre ($i = 0$) e, para cada torre seguinte, substituímos cada segmento ativo pela sua interseção com a borda do novo quadrado (que é um segmento ou até 2 pontos extremos!), mantendo no máximo $\Ocomp{1}$ segmentos ativos! Ao final, varremos os pontos inteiros com $u \equiv v \pmod 2$ nos segmentos sobreviventes!}
  {Se a primeira torre tiver $D_0 = 0$, os 4 segmentos degeneram em 1 único ponto. Ao intersectar um segmento horizontal $v = v_0, u \in [u_L, u_R]$ com a borda do quadrado $[U-D, U+D] \times [V-D, V+D]$: se $v_0 \notin [V-D, V+D]$, a interseção é vazia; se $v_0 \in \{V-D, V+D\}$ (toca a tampa superior/inferior do quadrado), a interseção é o sub-segmento $[ \max(u_L, U-D), \min(u_R, U+D) ]$; se $V - D < v_0 < V + D$, o segmento só toca as paredes laterais nos até 2 pontos $u = U - D$ e $u = U + D$ (que viram segmentos degenerados de comprimento 0)!}
  {Tempo: $\Ocomp{N + K \log K}$ ($K$ = pontos de saída) $\mid$ Espaço: $\Ocomp{K}$}

\begin{lstlisting}
struct OrthoSeg { ll u1, v1, u2, v2; }; /* u1 == u2 (vertical) ou v1 == v2 (horizontal) */

vector<pll> solve_gps_manhattan_spheres(const vector<tuple<ll, ll, ll>>& towers) {
  auto get_square_sides = [](ll x, ll y, ll d) -> vector<OrthoSeg> {
    ll U = x + y, V = x - y;
    if (d == 0) return {{U, V, U, V}};
    return {
      {U - d, V - d, U + d, V - d}, /* Base horizontal (v = V - d) */
      {U - d, V + d, U + d, V + d}, /* Topo horizontal (v = V + d) */
      {U - d, V - d, U - d, V + d}, /* Parede esquerda (u = U - d) */
      {U + d, V - d, U + d, V + d} /* Parede direita  (u = U + d) */
    };
  };
  /* Intersecta dois segmentos ortogonais a e b (retorna 0 ou 1 OrthoSeg) */
  auto isect_two = [](OrthoSeg a, OrthoSeg b, vector<OrthoSeg>& out) {
    if (a.u1 > a.u2) swap(a.u1, a.u2);
    if (a.v1 > a.v2) swap(a.v1, a.v2);
    if (b.u1 > b.u2) swap(b.u1, b.u2);
    if (b.v1 > b.v2) swap(b.v1, b.v2);
    ll uL = max(a.u1, b.u1), uR = min(a.u2, b.u2);
    ll vL = max(a.v1, b.v1), vR = min(a.v2, b.v2);
    if (uL <= uR && vL <= vR) out.pb({uL, vL, uR, vR});
  };

  auto [x0, y0, d0] = towers[0];
  vector<OrthoSeg> active = get_square_sides(x0, y0, d0);

  rep(i, 1, sz(towers)) {
    auto [xi, yi, di] = towers[i];
    auto sides = get_square_sides(xi, yi, di);
    vector<OrthoSeg> nxt;
    for (const auto& s1 : active)
      for (const auto& s2 : sides) isect_two(s1, s2, nxt);
    /* Deduplica segmentos identicos para manter sz(active) <= 8 */
    active.clear();
    for (const auto& s : nxt) {
      bool dup = false;
      for (const auto& ex : active)
        if (s.u1 == ex.u1 && s.v1 == ex.v1 && s.u2 == ex.u2 && s.v2 == ex.v2) { dup = true; break; }
      if (!dup) active.pb(s);
    }
    if (active.empty()) break;
  }
  /* Extrai todos os pontos inteiros (x, y) dos segmentos finais (garantido conjunto finito pelo problema) */
  vector<pll> ans;
  for (auto [u1, v1, u2, v2] : active) {
    for (ll u = u1; u <= u2; ++u) {
      for (ll v = v1; v <= v2; ++v) {
        if (((u ^ v) & 1) == 0) ans.eb((u + v) / 2, (u - v) / 2);
      }
    }
  }
  sort(all(ans));
  ans.erase(unique(all(ans)), ans.end());
  return ans;
}
\end{lstlisting}

\subsubsection{Interseção Reta-Parábola e Cruzamento de Cordas via BIT}
\algobox
  {Dadas $N$ retas distintas $y = M_i x + C_i$ e a região aberta $y > A x^2 + B$ ($A > 0$), contar quantos pares de retas se cruzam estritamente dentro da parábola.}
  {Uma reta $y = M_i x + C_i$ entra na região aberta $y > A x^2 + B$ sse a equação $A x^2 - M_i x + (B - C_i) = 0$ possui duas raízes reais distintas, ou seja, $\Delta_i = M_i^2 - 4A(B - C_i) > 0$! Nesse caso, a interseção da reta $i$ com o interior da parábola é uma \textbf{corda aberta} com abscissas extremas $(l_i, r_i) = \bigl(\frac{M_i - \sqrt{\Delta_i}}{2A}, \frac{M_i + \sqrt{\Delta_i}}{2A}\bigr)$ na fronteira da parábola!
  Como a parábola é estritamente convexa, duas cordas $(l_i, r_i)$ e $(l_j, r_j)$ se cruzam estritamente no interior da região parabólica sse seus extremos se alternam estritamente ao longo da fronteira da parábola: $l_i < l_j < r_i < r_j$ ou $l_j < l_i < r_j < r_i$!
  Isso reduz o problema a \textbf{contar pares de intervalos $(l_i, r_i)$ que se entrelaçam estritamente}, resolvido em $\Ocomp{N \log N}$ com compressão de coordenadas + Fenwick Tree (BIT)!}
  {Atenção à precisão e empates na fronteira: duas retas que se cruzam \emph{exatamente sobre} a parábola compartilham um extremo $l_i = l_j$, $r_i = r_j$ ou $r_i = l_j$, o que NÃO deve ser contado! Ordenando por $l_i$ crescente (e em empate de $l_i$ por $r_i$ decrescente) e usando tolerância $\eps = 10^{-11}$ em \texttt{long double} na compressão de coordenadas de $r_i$, excluímos todos os encontros na fronteira!}
  {Tempo: $\Ocomp{N \log N}$ $\mid$ Espaço: $\Ocomp{N}$}

\begin{lstlisting}
ll count_parabola_chord_intersections(ll A, ll B, const vector<pll>& lines) {
  constexpr ld EPS_P = 1e-11L;
  vector<pair<ld, ld>> chords;
  chords.reserve(sz(lines));
  for (auto [M, C] : lines) {
    ll delta = M * M - 4LL * A * (B - C);
    if (delta <= 0) continue; /* Nao entra no interior da parabola */
    ld sq = sqrtl((ld)delta);
    ld l = ((ld)M - sq) / (2.0L * A);
    ld r = ((ld)M + sq) / (2.0L * A);
    chords.eb(l, r);
  }
  int m = sz(chords);
  if (m < 2) return 0;

  /* Coleta e unifica coordenadas de r_i (e l_i) com tolerancia EPS_P */
  vector<ld> coords;
  coords.reserve(2 * m);
  for (auto [l, r] : chords) { coords.pb(l); coords.pb(r); }
  sort(all(coords));
  vector<ld> uniq;
  for (ld x : coords) {
    if (uniq.empty() || x - uniq.back() > EPS_P) uniq.pb(x);
  }
  auto get_id = [&](ld x) -> int {
    auto it = lower_bound(all(uniq), x - EPS_P);
    return (int)(it - uniq.begin());
  };

  vector<pii> segs(m);
  rep(i, 0, m) segs[i] = {get_id(chords[i].fst), get_id(chords[i].snd)};
  /* Ordena por L crescente; se L empatar (encontro na borda esquerda), ordena por R DECRESCENTE
   * para que o primeiro nao conte o segundo em query(L+1, R-1)! */
  sort(all(segs), [](const pii& a, const pii& b) {
    return a.fst != b.fst ? a.fst < b.fst : a.snd > b.snd;
  });

  int K = sz(uniq);
  Fenwick<int> bit(K);
  ll ans = 0;
  /* Agrupa cordas com mesmo L para nao inserir R antes de consultar o mesmo L */
  for (int i = 0; i < m; ) {
    int j = i;
    while (j < m && segs[j].fst == segs[i].fst) ++j;
    rep(k, i, j) {
      auto [L, R] = segs[k];
      if (L + 1 <= R - 1) ans += bit.query(L + 1, R - 1);
    }
    rep(k, i, j) bit.add(segs[k].snd, 1);
    i = j;
  }
  return ans;
}
\end{lstlisting}

\subsubsection{Busca Binária Paralela + CHT e Filtragem no Espaço-Tempo}
\algobox
  {(1) Dados $T$ pontos $(X_i, Y_i)$ e $P$ retas $y = A_j x + B_j$, achar para cada ponto a primeira reta $j$ com $Y_i < A_j X_i + B_j$; (2) Filtrar intervalos verticais $(U_j, [A_j, B_j])$ não cobertos por semirretas $x(t) = (t - T_i)S_i$ e responder queries $(T', S')$.}
  {\textbf{1. PBS + CHT:} Um ponto $(X_i, Y_i)$ é capturado no turno $\le m$ sse $\max_{1 \le j \le m} (A_j X_i + B_j) > Y_i$. Com Busca Binária Paralela (PBS) + CHT incremental, resolvemos todas as camadas em $\Ocomp{(T + P) \log P}$. \textbf{2. Janela Espaço-Tempo:} Ordenamos os segmentos por $U_j$ e filtramos os descobertos (parando em $T_i > U_j$). Para cada query $(T', S')$, restringimos a busca por busca binária à janela $U_j \in [T', T' + X_{\max} / S']$.}
  {Ordenar os pontos por $X_i$ antes da PBS melhora a localidade de cache no CHT.}
  {PBS + CHT: $\Ocomp{(T + P) \log P}$ $\mid$ Espaço: $\Ocomp{T + P}$}

\begin{lstlisting}
/* 1. Primeira reta estritamente acima de cada ponto em O((T + P) log P) via Busca Binaria Paralela + CHT Estatico */
vector<vi> solve_board_game_pbs(const vector<pll>& tokens, const vector<pll>& cards) {
  int T = sz(tokens), P = sz(cards);
  /* Ordena tokens por X crescente para consultar o CHT com ponteiro monotono em O(1)! */
  vi tok_ord(T); iota(all(tok_ord), 0);
  sort(all(tok_ord), [&](int i, int j) { return tokens[i].fst < tokens[j].fst; });

  vi lo(T, 0), hi(T, P - 1), ans(T, P);
  vector<vi> buckets(P);

  /* Stack para CHT estatico (inclinacao k crescente, consulta maxima) */
  struct CLine { ll k, m; };
  auto bad = [](const CLine& l1, const CLine& l2, const CLine& l3) {
    /* Intersecao de l1 e l2 >= intersecao de l2 e l3 <=> (l2.m - l1.m)*(l3.k - l2.k) >= (l3.m - l2.m)*(l2.k - l1.k) */
    return (i128)(l2.m - l1.m) * (l3.k - l2.k) >= (i128)(l3.m - l2.m) * (l2.k - l1.k);
  };

  while (true) {
    bool any = false;
    rep(i, 0, P) buckets[i].clear();
    for (int id : tok_ord) if (lo[id] <= hi[id]) {
      buckets[(lo[id] + hi[id]) >> 1].pb(id);
      any = true;
    }
    if (!any) break;

    /* Podemos usar LiChao / LineContainer incremental de 0 ate P-1 em uma unica varredura! */
    LineContainer cht;
    rep(mid, 0, P) {
      cht.add(cards[mid].fst, cards[mid].snd);
      for (int id : buckets[mid]) {
        auto [x, y] = tokens[id];
        if (cht.query(x) > y) { ans[id] = mid; hi[id] = mid - 1; }
        else                  { lo[id] = mid + 1; }
      }
    }
  }
  vector<vi> player_tokens(P);
  rep(id, 0, T) if (ans[id] < P) player_tokens[ans[id]].pb(id + 1); /* 1-based token IDs */
  return player_tokens;
}

/* 2. Filtragem de semirretas espaco-tempo contra segmentos verticais
 * Busca binaria na janela temporal ativa U in [t, t + MAX_COORD / s] */
vi solve_joining_marathon(const vector<pll>& runners,
                          const vector<tuple<ll, ll, ll>>& photos,
                          const vector<pll>& queries) {
  int P = sz(photos), R = sz(runners), Q = sz(queries);
  /* Ordena fotos por tempo U crescente para podar U < T e U > T + MAX_X / S via binary search! */
  struct Photo { int u, a, b; };
  vector<Photo> ph(P);
  rep(i, 0, P) {
    auto [u, a, b] = photos[i];
    ph[i] = {(int)u, (int)a, (int)b};
  }
  sort(all(ph), [](const Photo& x, const Photo& y) { return x.u < y.u; });

  /* Filtra apenas as fotos que sao LIXO (nenhum dos R corredores iniciais aparece nelas)
   * Ordena corredores por T crescente para parar assim que T_i > U */
  vector<pll> run_sorted = runners;
  sort(all(run_sorted));
  vector<Photo> trash;
  trash.reserve(P);
  for (const auto& p : ph) {
    bool covered = false;
    for (auto [t, s] : run_sorted) {
      if (t > p.u) break;
      ll pos = (p.u - t) * s;
      if (pos >= p.a && pos <= p.b) { covered = true; break; }
    }
    if (!covered) trash.pb(p);
  }
  int M = sz(trash);
  vi trash_u(M);
  rep(i, 0, M) trash_u[i] = trash[i].u;

  vi ans(Q, M);
  static constexpr ll MAX_COORD = 1'000'000'000LL;
  rep(q, 0, Q) {
    ll t = queries[q].fst, s = queries[q].snd;
    /* Uma foto so pode ser salva por (t, s) se t <= U <= t + MAX_COORD / s */
    ll max_u = t + MAX_COORD / s;
    int l = (int)(lower_bound(all(trash_u), (int)min(t, (ll)INT_MAX)) - trash_u.begin());
    int r = (int)(upper_bound(all(trash_u), (int)min(max_u, (ll)INT_MAX)) - trash_u.begin());
    int saved = 0;
    for (int i = l; i < r; ++i) {
      ll pos = (trash[i].u - t) * s;
      if (pos >= trash[i].a && pos <= trash[i].b) ++saved;
    }
    ans[q] = M - saved;
  }
  return ans;
}
\end{lstlisting}

% ----------------------------------------------------------------------------
% SUBTEMA 5.3: SWEEP RADIAL, SEMIPLANOS, CIRCULOS E MINKOWSKI
% ----------------------------------------------------------------------------
\subsection{Sweep Radial, Semiplanos, Círculos e Soma de Minkowski}

\subsubsection{Ordenação Polar Exata e Interseção de Intervalos Angulares}
\algobox
  {(1) Ordenar pontos/vetores por ângulo polar em $[-\pi, \pi)$ usando apenas inteiros (\texttt{half} + \texttt{cross}); (2) Para semiplanos de direção, decidir se existe uma direção de gravidade/impressão $\vec{g}$ tal que $\vec{g} \cdot P_i \ge \vec{g} \cdot P_j$ para todo par com $B_i > B_j$; (3) Para o maior quadrado vazio, achar o maior quadrado centrado em $(0,0)$ que não contém nenhum ponto no seu interior estrito.}
  {\textbf{1. Semiplanos de Direção:} Para cada par $(i, j)$ com $B_i > B_j$, o vetor direção $\vec{g}$ deve satisfazer $\vec{g} \cdot (P_i - P_j) \ge 0$, o que restringe o ângulo de $\vec{g}$ a um semicírculo fechado de comprimento $\pi$ centrado na direção de $\vec{v}_{ij} = P_i - P_j$! Existe $\vec{g}$ sse a interseção de todos esses semicírculos não for vazia, o que equivale a: todos os vetores $\vec{v}_{ij}$ cabem em um semicírculo fechado de abertura $\le \pi$ (basta ordenar os vetores $\vec{v}_{ij}$ radialmente e checar se existe um gap angular $\ge \pi$ entre dois vetores consecutivos, sem contar vetores exatamente opostos $\vec{v}$ e $-\vec{v}$ a menos que haja direção perpendicular comum!).
  \textbf{2. Maior Quadrado Vazio:} Como o quadrado centrado na origem tem simetria de $90^\circ$ ($\pi/2$), podemos rotacionar todos os $N$ pontos para o primeiro quadrante (ou testar o ângulo de inclinação $\theta \in [0, \pi/2)$ do quadrado)! Para um semilado $L$ fixo (distância da origem aos 4 lados do quadrado), um ponto $P_i$ a distância $r_i = |P_i|$ da origem impede os ângulos $\theta$ nos quais $|x'_i| < L$ e $|y'_i| < L$:
  $\bullet$ Se $r_i < L$: o ponto $P_i$ está dentro do círculo inscrito no quadrado, logo fica estritamente dentro do quadrado para \textbf{qualquer} ângulo ($L$ é impossível!).
  $\bullet$ Se $r_i \ge \sqrt{2} L$: o ponto está fora do círculo circunscrito ao quadrado, logo nunca entra no quadrado!
  $\bullet$ Se $L \le r_i < \sqrt{2} L$: seja $\alpha_i = \text{atan2}(y_i, x_i) \bmod (\pi/2) \in [0, \pi/2)$ e $\delta_i = \arccos(L / r_i) \in [0, \pi/4)$. Na rotação por $\theta$, o ponto entra no quadrado sse sua distância angular aos eixos fica entre $\delta_i$ e $\frac{\pi}{2} - \delta_i$, proibindo o intervalo aberto $(\alpha_i + \delta_i - \frac{\pi}{2}, \alpha_i - \delta_i) \pmod{\pi/2}$! Fazemos Busca Binária em $L$ + União de Intervalos em $[0, \pi/2)$!}
  {Nos semiplanos de direção, com implementação por ângulos ou produto vetorial: ordenar as direções únicas de $\vec{v}_{ij} = P_i - P_j$ radialmente; existe um semiplano contendo todos os $\vec{v}_{ij}$ sse há um salto angular $\ge \pi$ entre duas direções consecutivas $d_k$ e $d_{k+1}$ (isto é, $d_k \times d_{k+1} < 0$, ou $d_k \times d_{k+1} == 0$ com $d_k \cdot d_{k+1} < 0$ desde que todos os vetores estejam no mesmo semiplano fechado!).}
  {Semiplanos: $\Ocomp{N^2 \log N}$ $\mid$ Quadrado Vazio: $\Ocomp{N \log N \cdot \text{iter}}$}

\begin{lstlisting}
/* Predicado de ordenacao polar exata em [-pi, pi) sem ponto flutuante */
int half_plane(const Pll& p) { return p.y < 0 || (p.y == 0 && p.x < 0); }
bool polar_cmp(const Pll& a, const Pll& b) {
  int ha = half_plane(a), hb = half_plane(b);
  if (ha != hb) return ha < hb;
  return a.cross(b) > 0;
}

/* 1. Verifica existencia de semiplano comum de orientacao em O(N^2 log N) com aritmetica 100% inteira */
bool can_rotate_dazzling_stars(const vector<tuple<ll, ll, int>>& stars) {
  int n = sz(stars);
  vector<Pll> dirs;
  rep(i, 0, n) rep(j, 0, n) {
    if (get<2>(stars[i]) > get<2>(stars[j])) {
      ll dx = get<0>(stars[i]) - get<0>(stars[j]);
      ll dy = get<1>(stars[i]) - get<1>(stars[j]);
      ll g = std::gcd(abs(dx), abs(dy));
      dirs.pb({dx / g, dy / g});
    }
  }
  if (dirs.empty()) return true;
  sort(all(dirs), polar_cmp);
  dirs.erase(unique(all(dirs)), dirs.end());
  int m = sz(dirs);
  if (m == 1) return true;
  /* Existe um semiplano fechado que contem todos os vetores (com produto escalar >= 0)
   * sse existe um par adjacente (dirs[i], dirs[(i+1)%m]) com giro >= 180 graus (CCW <= 0)! */
  rep(i, 0, m) {
    Pll a = dirs[i], b = dirs[(i + 1) % m];
    ll cr = a.cross(b), dt = a.dot(b);
    if (cr < 0 || (cr == 0 && dt < 0)) return true;
  }
  return false;
}

/* 2. Maior quadrado centrado na origem sem pontos internos via Busca Binaria + Sweep Angular mod pi/2 */
ld max_perfect_house_perimeter(const vector<Pll>& plants) {
  int n = sz(plants);
  ld Period = PI / 2.0L;
  vector<ld> r(n), alpha(n);
  ld min_r = 1e30L;
  rep(i, 0, n) {
    r[i] = plants[i].dist();
    chmin(min_r, r[i]);
    ld a = fmodl(plants[i].angle(), Period);
    if (a < 0) a += Period;
    alpha[i] = a;
  }
  /* Para o semilado L (perimetro = 8 * L):
   * Sabemos que min_r / sqrt(2) <= L_opt <= min_r! */
  auto can_place = [&](ld L) -> bool {
    vector<pair<ld, ld>> bad;
    bad.reserve(2 * n);
    rep(i, 0, n) {
      if (r[i] < L) return false;
      if (r[i] >= L * sqrtl(2.0L)) continue;
      /* Ponto entra no interior estrito (-L, L)^2 quando phi = (alpha[i] - theta) mod (pi/2)
       * satisfaz r[i]*cos(phi) < L E r[i]*sin(phi) < L <=> delta < phi < pi/2 - delta! */
      ld delta = acosl(clamp(L / r[i], -1.0L, 1.0L));
      ld len = Period - 2.0L * delta; /* Largura do intervalo proibido */
      if (len <= 1e-13L) continue;
      /* phi in (delta, pi/2 - delta) <=> theta in (alpha[i] - (pi/2 - delta), alpha[i] - delta) mod Period */
      ld l = alpha[i] - (Period - delta), h = alpha[i] - delta;
      while (l < 0) { l += Period; h += Period; }
      while (l >= Period) { l -= Period; h -= Period; }
      if (h <= Period) bad.eb(l, h);
      else { bad.eb(l, Period); bad.eb(0.0L, h - Period); }
    }
    return union_of_intervals_length(move(bad)) < Period - 1e-12L;
  };

  ld lo = min_r / sqrtl(2.0L), hi = min_r;
  rep(iter, 0, 60) {
    ld mid = (lo + hi) * 0.5L;
    if (can_place(mid)) lo = mid;
    else                hi = mid;
  }
  return 8.0L * lo;
}
\end{lstlisting}

\subsubsection{DP Radial de Polígonos Visíveis}
\algobox
  {Dados a casa $H$, o portão $G$ e $N \le 300$ torres, contar quantos polígonos fechados têm $G$ e um subconjunto de torres como vértices tais que: (1) $H$ é estritamente interno ao polígono; (2) todo o muro é visível a partir de $H$ (polígono estrelado que dá exatamente 1 volta em torno de $H$); (3) todos os ângulos internos nas torres são $< 180^\circ$ ($\bmod 10^9+7$).}
  {Transladamos a casa $H$ para a origem $(0, 0)$. Para que todo o muro seja visível a partir de $H$ e encerre $H$ estritamente, os vértices escolhidos $G = v_0, v_1, \dots, v_k, G$ devem progredir \textbf{estritamente no sentido anti-horário (CCW)} em torno de $H$, com cada aresta $(v_i, v_{i+1})$ dando um giro estritamente entre $0^\circ$ e $180^\circ$ vista de $H$, isto é, $\text{cross}(H, v_i, v_{i+1}) > 0$!
  E a condição de que o ângulo interno na torre $v_i$ seja $< 180^\circ$ equivale a dizer que a curva $v_{i-1} \to v_i \to v_{i+1}$ também dobra estritamente para a esquerda (CCW): $\text{cross}(v_{i-1}, v_i, v_{i+1}) > 0$!
  Logo: ordenamos todas as $N$ torres em sentido anti-horário ao redor de $H$ começando a partir do raio $H \to G$ (descartando pontos colineares com $H$ e $G$ que dariam giro $0^\circ$ ou $180^\circ$; se duas torres têm o mesmo raio a partir de $H$, não podemos usar ambas consecutivamente!). Definimos $dp[i][j]$ = número de cadeias válidas $G \to \dots \to i \to j$: para transitar de $(i, j)$ para $k > j$ (com $\text{cross}(H, j, k) > 0$), exigimos $\text{cross}(i, j, k) > 0$! Ao ordenar para cada $j$ os anteriores $i$ e os próximos $k$ radialmente em torno de $j$, dois ponteiros + soma de prefixos reduzem a transição de $\Ocomp{N^3}$ para $\Ocomp{N^2 \log N}$!}
  {Para que dois muros sejam considerados distintos sse diferem no conjunto de torres usadas, basta notar que todo conjunto de torres que forma um polígono visível em torno de $H$ possui uma \textbf{única} ordem cíclica anti-horária ao redor de $H$, exceto se duas torres estiverem exatamente no mesmo raio a partir de $H$ (mas como $\text{cross}(H, v_i, v_{i+1}) > 0$ é obrigatório para a visibilidade estrita, nunca podemos escolher duas torres no mesmo raio!).}
  {Tempo: $\Ocomp{N^3}$ (ou $\Ocomp{N^2 \log N}$) $\mid$ Espaço: $\Ocomp{N^2}$}

\begin{lstlisting}
ll count_impenetrable_walls(Pll H, Pll G, vector<Pll> towers) {
  constexpr ll MOD = 1'000'000'007;
  G = G - H;
  vector<Pll> pts;
  for (auto p : towers) {
    p = p - H;
    /* Torres colineares na mesma semirreta de G nao podem ser usadas com visibilidade estrita */
    if (G.cross(p) == 0 && G.dot(p) > 0) continue;
    pts.pb(p);
  }
  /* Ordena todas as torres no sentido anti-horario (CCW) comecando logo apos o vetor G */
  auto ccw_from_G = [&](const Pll& a, const Pll& b) {
    ll ca = G.cross(a), cb = G.cross(b);
    bool ha = (ca > 0 || (ca == 0 && G.dot(a) < 0)); /* 1 se a esta em (0, pi] a partir de G */
    bool hb = (cb > 0 || (cb == 0 && G.dot(b) < 0));
    if (ha != hb) return ha > hb;
    return a.cross(b) > 0;
  };
  sort(all(pts), ccw_from_G);

  int n = sz(pts);
  pts.insert(pts.begin(), G);
  /* dp[u][v] = caminhos 0 -> ... -> u -> v (1 <= v <= n, 0 <= u < v) */
  vector<vll> dp(n + 1, vll(n + 1, 0));
  rep(v, 1, n + 1) {
    if (pts[0].cross(pts[v]) > 0) dp[0][v] = 1;
  }
  ll ans = 0;
  rep(v, 1, n + 1) {
    rep(u, 0, v) if (dp[u][v] > 0) {
      ll ways = dp[u][v];
      /* Tenta fechar o poligono voltando para G (indice 0): precisa pts[v].cross(G) > 0
       * e angulo interno em v < 180 graus (ccw(pts[u], pts[v], G) > 0) */
      if (u != 0 && pts[v].cross(pts[0]) > 0 && ccw(pts[u], pts[v], pts[0]) > 0) {
        ans = (ans + ways) % MOD;
      }
      /* Transita para a proxima torre w > v */
      rep(w, v + 1, n + 1) {
        if (pts[v].cross(pts[w]) > 0 && ccw(pts[u], pts[v], pts[w]) > 0) {
          dp[v][w] += ways;
          if (dp[v][w] >= MOD) dp[v][w] -= MOD;
        }
      }
    }
  }
  return ans;
}
\end{lstlisting}

\subsubsection{Interseção de Semiplanos em $\Ocomp{N \log N}$ e Reta-Círculo}
\algobox
  {(1) Calcular o polígono convexo resultante da interseção de $N$ semiplanos $\vec{n}_i \cdot X \le c_i$ (região à esquerda da reta orientada $P_i \to Q_i$) em $\Ocomp{N \log N}$; (2) Para bloquear a origem em raio $D$, encontrar o menor prefixo de retas que impede a origem $(0,0)$ de alcançar qualquer ponto na circunferência de raio $D$.}
  {\textbf{Bloqueio da Origem ao Raio $D$:}
  \textbf{Método 1 (Busca Binária + Interseção de Semiplanos):} Cada reta que não passa por $(0,0)$ define um semiplano contendo $(0,0)$. O urso consegue chegar a algum ponto a distância $D$ da origem sse a interseção dos primeiros $k$ semiplanos contém algum ponto com $|P| \ge D$, o que pela convexidade equivale a: a interseção dos $k$ semiplanos é ilimitada OU possui algum vértice $V$ com $|V| \ge D$! (Para garantir que a região é sempre limitada na busca binária, adicionamos uma caixa gigante $[-2D, 2D]^2$!).
  \textbf{Método 2 ($\Ocomp{N \log N}$ Direto via Arcos ou Dual):} Basta usar Busca Binária com a struct \texttt{HalfplaneIntersection} pronta abaixo!}
  {Na \texttt{HalfplaneIntersection}, adicionar primeiro uma \emph{Bounding Box} de lado $M = 2 \cdot D + 10$ transforma qualquer região aberta em um polígono convexo fechado: então basta verificar se algum vértice desse polígono satisfaz $x^2 + y^2 \ge D^2$!}
  {Interseção de Semiplanos: $\Ocomp{N \log N}$ $\mid$ Com Busca Binária: $\Ocomp{N \log^2 N}$}

\begin{lstlisting}
struct Halfplane {
  Pld p, pq; /* Reta passando por 'p' com direcao 'pq' (semiplano a ESQUERDA / CCW de pq) */
  ld ang;
  Halfplane() {}
  Halfplane(Pld a, Pld b) : p(a), pq(b - a), ang(pq.angle()) {}

  bool out(const Pld& r) const { return pq.cross(r - p) < -EPS; }
  bool operator<(const Halfplane& o) const {
    if (fabsl(ang - o.ang) < EPS) return pq.cross(o.p - p) < 0;
    return ang < o.ang;
  }
  friend Pld isect(const Halfplane& s, const Halfplane& t) {
    ld alpha = (t.p - s.p).cross(t.pq) / s.pq.cross(t.pq);
    return s.p + s.pq * alpha;
  }
};

/* Retorna os vertices (em CCW) da intersecao dos semiplanos (vazio se inviavel) */
vector<Pld> halfplane_intersection(vector<Halfplane> H) {
  sort(all(H));
  deque<Halfplane> dq;
  rep(i, 0, sz(H)) {
    if (i > 0 && fabsl(H[i].ang - H[i - 1].ang) < EPS) continue;
    while (sz(dq) >= 2 && H[i].out(isect(dq.back(), dq[sz(dq) - 2]))) dq.pop_back();
    while (sz(dq) >= 2 && H[i].out(isect(dq[0], dq[1]))) dq.pop_front();
    dq.pb(H[i]);
  }
  while (sz(dq) >= 3 && dq.back().out(isect(dq[0], dq[1]))) dq.pop_front();
  while (sz(dq) >= 3 && dq[0].out(isect(dq.back(), dq[sz(dq) - 2]))) dq.pop_back();
  if (sz(dq) < 3) return {};
  vector<Pld> poly(sz(dq));
  rep(i, 0, sz(dq)) poly[i] = isect(dq[i], dq[(i + 1) % sz(dq)]);
  return poly;
}

/* Menor prefixo de semiplanos que confina a origem no raio D: Retorna 1-based indice da predicao ou -1 se "*" */
int earliest_hazard_prediction(ld D, const vector<pair<Pld, Pld>>& lines) {
  int n = sz(lines);
  ld B = 2.0L * D + 100.0L;
  vector<Halfplane> box = {
    Halfplane({-B, -B}, {B, -B}), Halfplane({B, -B}, {B, B}),
    Halfplane({B, B}, {-B, B}),   Halfplane({-B, B}, {-B, -B})
  };
  vector<Halfplane> all_hp(n);
  rep(i, 0, n) {
    auto [a, b] = lines[i];
    /* Orienta a -> b de modo que a origem (0,0) esteja estritamente a esquerda (CCW) */
    if ((b - a).cross(Pld{0, 0} - a) < 0) swap(a, b);
    all_hp[i] = Halfplane(a, b);
  }
  auto can_reach_D = [&](int k) -> bool {
    vector<Halfplane> cur = box;
    cur.insert(cur.end(), all_hp.begin(), all_hp.begin() + k);
    auto poly = halfplane_intersection(move(cur));
    ld D2 = D * D;
    for (const auto& v : poly) if (v.dist2() >= D2 - 1e-7L) return true;
    return false;
  };
  if (can_reach_D(n)) return -1; /* Nunca fica bloqueado ("*") */
  int lo = 1, hi = n, ans = n;
  while (lo <= hi) {
    int mid = (lo + hi) >> 1;
    if (!can_reach_D(mid)) { ans = mid; hi = mid - 1; }
    else                   { lo = mid + 1; }
  }
  return ans;
}
\end{lstlisting}

\subsubsection{Sweep Radial de Todas as $\binom{N}{2}$ Retas — Rotating Line}
\algobox
  {Dados $N \le 400$ pontos sem 3 colineares e área mínima $S$, contar quantos quadriláteros simples (convexos ou côncavos!) possuem vértices na entrada e área $\ge S$.}
  {\textbf{1. Decomposição de Quadriláteros Simples (Convexos vs. Côncavos):}
  $\bullet$ Todo conjunto de 4 pontos em posição convexa forma \textbf{exatamente 1} quadrilátero simples, que possui \textbf{2 diagonais internas} (cada diagonal separa 2 vértices de um lado e 2 do outro).
  $\bullet$ Todo conjunto de 4 pontos em posição côncava (1 ponto interno ao triângulo dos outros 3) forma \textbf{exatamente 3} quadriláteros simples (côncavos), e cada um desses 3 quadriláteros possui \textbf{exatamente 1 diagonal interna} (que separa os outros 2 vértices em lados opostos da reta suporte da diagonal) e 1 diagonal externa (onde os outros 2 vértices estão do \textbf{mesmo lado} da reta suporte e o triângulo maior contém o triângulo menor)!
  \textbf{2. Contagem Unificada pelas Diagonais via Rotating Line Sweep em $\Ocomp{N^2 \log N}$:}
  Mantenha os $N$ pontos ordenados pela projeção perpendicular a uma reta que gira $180^\circ$ passando por todos os $\binom{N}{2}$ pares de pontos $(A, B)$:
  $\bullet$ \textbf{Lados Opostos de $AB$ (Diagonal Interna):} Os pontos abaixo de $AB$ estão em \texttt{ord[0 .. posA - 1]} (com área do triângulo $\triangle PAB$ decrescente à medida que o índice se aproxima de $A$), e os pontos acima estão em \texttt{ord[posB + 1 .. N - 1]} (com área crescente)! Para cada ponto $P$ de um lado com $2 \cdot \text{Area}(\triangle PAB) = a_P$, contamos por Busca Binária (ou Two Pointers!) quantos pontos $Q$ do lado oposto satisfazem $a_P + a_Q \ge 2S$! Note que essa soma sobre todas as $\binom{N}{2}$ diagonais internas conta cada quadrilátero \textbf{convexo} válido \textbf{2 vezes} (pelas suas 2 diagonais) e cada quadrilátero \textbf{côncavo} válido \textbf{1 vez} (pela sua única diagonal interna)!
  $\bullet$ \textbf{Separando Convexos e Côncavos:} Um quadrilátero côncavo tem um ponto interno $D$ dentro de um triângulo $\triangle ABC$, formando 3 quadriláteros simples côncavos de áreas: $\text{Area}(\triangle ABC) - \text{Area}(\triangle ABD)$, $\text{Area}(\triangle ABC) - \text{Area}(\triangle BCD)$ e $\text{Area}(\triangle ABC) - \text{Area}(\triangle CAD)$! Note que no quadrilátero côncavo de diagonal externa $AB$, a área é $|\text{Area}(\triangle CAB) - \text{Area}(\triangle DAB)|$ onde $D$ está dentro de $\triangle CAB$! Ou mais direto: na contagem pelas diagonais internas, um par $(P, Q)$ em lados opostos de $AB$ forma um quadrilátero convexo sse o segmento $PQ$ cruza o segmento $AB$, e forma um quadrilátero côncavo sse $PQ$ \textbf{não} cruza $AB$!
  Ou ainda mais elegante: na soma pelas diagonais que se cruzam, toda diagonal que se cruza vem de um quádruplo convexo ($1 \times$ por quádruplo convexo), e para os côncavos fixamos o segmento $AB$ e contamos com a reta rotativa!}
  {Na verdade, para $N \le 400$, podemos contar exatamente: (1) os pares $(P, Q)$ em lados opostos de $AB$ com $a_P + a_Q \ge 2S$ para todos os pares $\{A, B\}$ (soma $T_{\text{opp}} = 2 \cdot Q_{\text{conv}} + Q_{\text{conc}}$), e (2) os pares $(C, D)$ do \textbf{mesmo lado} de $AB$ onde $D \in \triangle CAB$ e $a_C - a_D \ge 2S$: como todo quadrilátero côncavo tem exatamente 1 diagonal interna (contada em $T_{\text{opp}}$) e 1 diagonal externa (que é a aresta $AB$ oposta ao vértice reflexo $D \in \triangle CAB$, com área $a_C - a_D$!), a quantidade de quadriláteros côncavos válidos é \textbf{exatamente} $Q_{\text{conc}} = T_{\text{ext}}$! E então o total de quadriláteros simples válidos é $\frac{T_{\text{opp}} - Q_{\text{conc}}}{2} + Q_{\text{conc}} = \frac{T_{\text{opp}} + Q_{\text{conc}}}{2}$!}
  {Tempo: $\Ocomp{N^2 \log N}$ (ou $\Ocomp{N^3}$ com bitset) $\mid$ Espaço: $\Ocomp{N^2}$}

\begin{lstlisting}
/* Contagem de quadrilateros simples com area >= S em O(N^2 log N + N^3 / 64) */
ll count_simple_quadrilaterals_area(ll S, vector<Pll> pts) {
  int n = sz(pts);
  ll twoS = 2 * S; /* Trabalhamos com o dobro da area (inteiro exato!) */

  /* 1. Ordena pontos por (x, y) e pre-calcula bitset left_of[i][j] dos pontos a esquerda de (i -> j) */
  sort(all(pts));
  static bitset<405> left_of[405][405];
  rep(i, 0, n) rep(j, 0, n) left_of[i][j].reset();
  rep(i, 0, n) rep(j, 0, n) if (i != j) {
    rep(k, 0, n) if (k != i && k != j && pts[i].cross(pts[j], pts[k]) > 0)
      left_of[i][j].set(k);
  }

  vi ord(n), pos(n);
  iota(all(ord), 0); iota(all(pos), 0);
  struct Ev { int i, j; Pll v; };
  vector<Ev> evs;
  evs.reserve(n * (n - 1) / 2);
  rep(i, 0, n) rep(j, i + 1, n) evs.pb({i, j, pts[j] - pts[i]});
  sort(all(evs), [](const Ev& a, const Ev& b) {
    return polar_cmp(a.v, b.v);
  });

  ll T_opp = 0;
  for (const auto& e : evs) {
    int u = e.i, v = e.j;
    int pu = pos[u], pv = pos[v];
    if (pu > pv) { swap(u, v); swap(pu, pv); }
    int q_ptr = n - 1;
    for (int p_idx = pu - 1; p_idx >= 0; --p_idx) {
      ll area_p = abs(pts[u].cross(pts[v], pts[ord[p_idx]]));
      while (q_ptr > pv && area_p + abs(pts[u].cross(pts[v], pts[ord[q_ptr]])) >= twoS) {
        --q_ptr;
      }
      T_opp += (n - 1 - q_ptr);
    }
    swap(ord[pu], ord[pv]);
    swap(pos[u], pos[v]);
  }

  /* 3. Conta Q_conc: para cada par ordenado (A, B) e C a esquerda de (A -> B), */
  ll Q_conc = 0;
  vi left_list; left_list.reserve(n);
  vi rank_in_ab(n, 0);
  vll areas_ab; areas_ab.reserve(n);

  rep(a, 0, n) rep(b, 0, n) if (a != b) {
    left_list.clear();
    rep(c, 0, n) if (left_of[a][b].test(c)) left_list.pb(c);
    int m = sz(left_list);
    if (m < 2) continue;
    sort(all(left_list), [&](int x, int y) {
      return pts[a].cross(pts[b], pts[x]) < pts[a].cross(pts[b], pts[y]);
    });
    areas_ab.resize(m);
    rep(idx, 0, m) {
      int c = left_list[idx];
      rank_in_ab[c] = idx;
      areas_ab[idx] = pts[a].cross(pts[b], pts[c]);
    }
    rep(idx_c, 1, m) {
      ll max_d_area = areas_ab[idx_c] - twoS;
      if (max_d_area <= 0) continue;
      int lim = (int)(upper_bound(areas_ab.begin(), areas_ab.begin() + idx_c, max_d_area) - areas_ab.begin());
      int c = left_list[idx_c];
      const auto& mask_bc = left_of[b][c];
      const auto& mask_ca = left_of[c][a];
      rep(idx_d, 0, lim) {
        int d = left_list[idx_d];
        if (mask_bc.test(d) && mask_ca.test(d)) ++Q_conc;
      }
    }
  }
  return (T_opp + Q_conc) / 2;
}
\end{lstlisting}

\subsubsection{Envoltória Convexa, Soma de Minkowski e Ponto em Polígono Convexo $\Ocomp{\log N}$}
\algobox
  {(1) Construir a Envoltória Convexa (Monotone Chain) em $\Ocomp{N \log N}$; (2) Calcular a \textbf{Soma de Minkowski} $P + Q = \{p + q : p \in P, q \in Q\}$ de dois polígonos convexos em $\Ocomp{|P| + |Q|}$; (3) Testar se um ponto pertence a um polígono convexo (incluindo a borda) em $\Ocomp{\log N}$; (4) Para pontos médios entre polígonos, dados dois polígonos convexos $P_1, P_2$, decidir para cada ponto $q$ se existem $p_1 \in P_1, p_2 \in P_2$ colineares com $q$ tais que um dos três é ponto médio dos outros dois!}
  {\textbf{Ponto Médio via Soma de Minkowski:} Três pontos $p_1 \in P_1, p_2 \in P_2, q$ têm um deles exatamente no meio dos outros dois sse uma das 3 equações vetoriais vale:
  (1) $q$ é o ponto médio de $p_1$ e $p_2 \iff 2q = p_1 + p_2 \iff 2q \in P_1 + P_2$;
  (2) $p_1$ é o ponto médio de $q$ e $p_2 \iff q = 2p_1 - p_2 \iff q \in 2P_1 + (-P_2)$;
  (3) $p_2$ é o ponto médio de $q$ e $p_1 \iff q = 2p_2 - p_1 \iff q \in 2P_2 + (-P_1)$!
  Como $P_1, P_2$ são polígonos convexos, os conjuntos $M_1 = P_1 + P_2$, $M_2 = 2P_1 - P_2$ e $M_3 = 2P_2 - P_1$ são \textbf{polígonos convexos} calculados em $\Ocomp{|P_1| + |P_2|}$ via Soma de Minkowski! Então para cada uma das $N \le 5 \cdot 10^5$ usinas menores $q$, basta testar em $\Ocomp{\log M}$ se $2q \in M_1$ OU $q \in M_2$ OU $q \in M_3$!}
  {Antes de somar dois polígonos em Minkowski, rotacionar ciclicamente os vértices de cada polígono para que o vértice de menor $(y, x)$ fique no índice $0$, e remover vértices colineares redundantes após a soma para que o teste binário de ponto em polígono convexo em $\Ocomp{\log M}$ funcione sem degenerescências!}
  {Minkowski: $\Ocomp{|P| + |Q|}$ $\mid$ Query Ponto em Polígono Convexo: $\Ocomp{\log |P|}$}

\begin{lstlisting}
/* Reordena poligono convexo CCW para comecar no vertice de menor (y, x) */
void reorder_polygon(vector<Pll>& P) {
  int pos = 0;
  rep(i, 1, sz(P)) {
    if (P[i].y < P[pos].y || (P[i].y == P[pos].y && P[i].x < P[pos].x)) pos = i;
  }
  rotate(P.begin(), P.begin() + pos, P.end());
}

/* Soma de Minkowski de dois poligonos convexos CCW em O(|P| + |Q|) */
vector<Pll> minkowski_sum(vector<Pll> P, vector<Pll> Q) {
  reorder_polygon(P);
  reorder_polygon(Q);
  int n = sz(P), m = sz(Q);
  P.pb(P[0]); P.pb(P[1]);
  Q.pb(Q[0]); Q.pb(Q[1]);
  vector<Pll> res;
  res.reserve(n + m);
  int i = 0, j = 0;
  while (i < n || j < m) {
    res.pb(P[i] + Q[j]);
    ll cr = (P[i + 1] - P[i]).cross(Q[j + 1] - Q[j]);
    if (cr >= 0 && i < n) ++i;
    if (cr <= 0 && j < m) ++j;
  }
  /* Remove vertices colineares internos nas arestas para permitir Point-in-Convex-Polygon O(log N) */
  vector<Pll> clean;
  for (const auto& pt : res) {
    while (sz(clean) >= 2 && ccw(clean[sz(clean) - 2], clean.back(), pt) == 0)
      clean.pop_back();
    clean.pb(pt);
  }
  if (sz(clean) >= 3 && ccw(clean[sz(clean) - 2], clean.back(), clean[0]) == 0)
    clean.pop_back();
  return clean;
}

/* Testa se o ponto 'pt' pertence ao poligono convexo CCW 'P' (borda inclusa) em O(log |P|)
 * Requer que 'P' nao tenha 3 vertices colineares! */
bool point_in_convex_polygon(const vector<Pll>& P, const Pll& pt) {
  int n = sz(P);
  if (n == 1) return pt == P[0];
  if (n == 2) return on_segment(P[0], P[1], pt);
  if (ccw(P[0], P[1], pt) < 0 || ccw(P[0], P[n - 1], pt) > 0) return false;
  if (ccw(P[0], P[1], pt) == 0) return on_segment(P[0], P[1], pt);
  if (ccw(P[0], P[n - 1], pt) == 0) return on_segment(P[0], P[n - 1], pt);

  int lo = 1, hi = n - 2, idx = 1;
  while (lo <= hi) {
    int mid = (lo + hi) >> 1;
    if (ccw(P[0], P[mid], pt) >= 0) { idx = mid; lo = mid + 1; }
    else                            { hi = mid - 1; }
  }
  return ccw(P[idx], P[idx + 1], pt) >= 0;
}

/* Ponto medio colinear entre dois poligonos convexos via Minkowski em O((M1 + M2) + N log(M1 + M2)) */
string solve_gravitational_wave_detector(const vector<Pll>& P1, const vector<Pll>& P2,
                                         const vector<Pll>& queries) {
  auto scale_poly = [](vector<Pll> P, ll k) {
    for (auto& p : P) p = p * k;
    if (k < 0) reverse(P.begin() + 1, P.end()); /* Preserva orientacao CCW ao multiplicar por < 0! */
    return P;
  };
  vector<Pll> M1 = minkowski_sum(P1, P2); /* P1 + P2 (testar 2*q) */
  vector<Pll> M2 = minkowski_sum(scale_poly(P1, 2), scale_poly(P2, -1)); /* 2*P1 - P2 (testar q) */
  vector<Pll> M3 = minkowski_sum(scale_poly(P2, 2), scale_poly(P1, -1)); /* 2*P2 - P1 (testar q) */

  string ans(sz(queries), 'N');
  rep(i, 0, sz(queries)) {
    Pll q = queries[i];
    if (point_in_convex_polygon(M1, q * 2) ||
        point_in_convex_polygon(M2, q) ||
        point_in_convex_polygon(M3, q)) {
      ans[i] = 'Y';
    }
  }
  return ans;
}
\end{lstlisting}


\subsubsection{Geometria de Círculos: Interseções, Tangentes e Área Comum}
\algobox
  {Interseção Reta-Círculo, Interseção Círculo-Círculo, Tangentes a partir de ponto externo e área de interseção de dois círculos.}
  {Projetar o centro $O$ na reta $AB$ dá o pé da perpendicular $P$; a semidistância da corda é $h = \sqrt{r^2 - |P - O|^2}$. Para dois círculos $(O_1, r_1)$ e $(O_2, r_2)$ a distância $d = |O_2 - O_1|$, a lei dos cossenos fornece $a = \frac{r_1^2 - r_2^2 + d^2}{2d}$ e $h = \sqrt{r_1^2 - a^2}$.}
  {Tratar círculos concêntricos ($d == 0$) antes de dividir por $d$, e usar \texttt{clamp} antes de \texttt{acosl}/\texttt{sqrtl}.}
  {Tempo: $\Ocomp{1}$ por operação}

\begin{lstlisting}
/* 1. Intersecao Reta (a, b) com Circulo (c, r): retorna 0, 1 ou 2 pontos */
vector<Pld> circle_line_isect(Pld c, ld r, Pld a, Pld b) {
  Pld ab = b - a, p = a + ab * ((c - a).dot(ab) / ab.dist2());
  ld s = (c - p).dist2(), h2 = r * r - s;
  if (h2 < -EPS) return {};
  if (h2 < EPS) return {p};
  Pld dir = ab * (sqrtl(max(0.0L, h2)) / ab.dist());
  return {p - dir, p + dir};
}

/* 2. Intersecao de dois Circulos (c1, r1) e (c2, r2): retorna 0, 1 ou 2 pontos */
vector<Pld> circle_circle_isect(Pld c1, ld r1, Pld c2, ld r2) {
  Pld d_vec = c2 - c1;
  ld d2 = d_vec.dist2(), d = sqrtl(d2);
  if (d < EPS || d > r1 + r2 + EPS || d < fabsl(r1 - r2) - EPS) return {};
  ld a = (r1 * r1 - r2 * r2 + d2) / (2.0L * d);
  ld h2 = r1 * r1 - a * a;
  Pld mid = c1 + d_vec * (a / d);
  if (h2 < EPS) return {mid};
  Pld perp = d_vec.perp() * (sqrtl(max(0.0L, h2)) / d);
  return {mid - perp, mid + perp};
}

/* 3. Area de intersecao de dois circulos (c1, r1) e (c2, r2) */
ld circle_circle_area(Pld c1, ld r1, Pld c2, ld r2) {
  ld d = (c2 - c1).dist();
  if (d >= r1 + r2) return 0.0L;
  if (d <= fabsl(r1 - r2)) { ld rm = min(r1, r2); return PI * rm * rm; }
  ld a1 = 2.0L * acosl(clamp((d * d + r1 * r1 - r2 * r2) / (2.0L * d * r1), -1.0L, 1.0L));
  ld a2 = 2.0L * acosl(clamp((d * d + r2 * r2 - r1 * r1) / (2.0L * d * r2), -1.0L, 1.0L));
  return 0.5L * r1 * r1 * (a1 - sinl(a1)) + 0.5L * r2 * r2 * (a2 - sinl(a2));
}

/* 4. Tangentes comuns a dois circulos (c1, r1) e (c2, r2) (para ponto-circulo, use r2 = 0) */
vector<pair<Pld, Pld>> circle_tangents(Pld c1, ld r1, Pld c2, ld r2, bool inner) {
  if (inner) r2 = -r2;
  Pld d = c2 - c1;
  ld dr = r1 - r2, d2 = d.dist2(), h2 = d2 - dr * dr;
  if (d2 < EPS || h2 < -EPS) return {};
  h2 = max(0.0L, h2);
  vector<pair<Pld, Pld>> out;
  for (ld sign : {-1.0L, 1.0L}) {
    Pld v = (d * dr + d.perp() * (sqrtl(h2) * sign)) / d2;
    out.eb(c1 + v * r1, c2 + v * r2);
    if (h2 < EPS) break;
  }
  return out;
}
\end{lstlisting}
